% Generated draft: requires mathematical and source review.
\documentclass[11pt]{article}
\usepackage[margin=25mm]{geometry}
\usepackage{fontspec}
\setmainfont{Latin Modern Roman}
\usepackage{amsmath,amssymb,mathtools,unicode-math}
\setmathfont{Latin Modern Math}
\usepackage{enumitem,needspace,xurl,hyperref}
\setlength{\parindent}{0pt}
\setlength{\parskip}{5pt}
\setlength{\emergencystretch}{3em}
\newcommand{\N}{\mathbb{N}}
\newcommand{\Z}{\mathbb{Z}}
\newcommand{\Q}{\mathbb{Q}}
\newcommand{\R}{\mathbb{R}}
\newcommand{\C}{\mathbb{C}}
\newcommand{\F}{\mathbb{F}}
\newcommand{\A}{\mathbb{A}}
\newcommand{\PP}{\mathbb{P}}
\newcommand{\E}{\mathbb{E}}
\newcommand{\eps}{\varepsilon}
\newcommand{\dd}{\,\mathrm{d}}
\begin{document}
\section*{UnsolvedMath ID: 30005455}
\textbf{Source problem code:} catalogue entry
\subsection{Definitions and mathematical statement}
The supplied record concerns supercritical reinforcement, with parameter \(\alpha>1\), on an underlying graph \(G\), with vertex firing rates \(p(v)\). It does not specify the reinforcement dynamics, admissible graphs, initial conditions, or construction of the surviving graph. Write \(H_\alpha\) provisionally for that surviving graph, and \(\mathcal C(H_\alpha)\) for its connected components. This notation does not supply the missing construction. A tree is a connected undirected graph containing no cycle. The source does not define a \emph{whisker}; that term must retain its source-specific meaning pending verification.

The intended problem is to improve the supplied sufficient bounds \(\alpha>2\) for the assertion
\[
\text{every }C\in\mathcal C(H_\alpha)\text{ is a tree},
\]
and \(\alpha>25\), under the additional assumption \(p(v)=1\) for every vertex \(v\), for the assertion
\[
\text{every }C\in\mathcal C(H_\alpha)\text{ is a whisker}.
\]
In particular, determine whether both assertions hold throughout \(\alpha>4/3\), under their respective original hypotheses, and whether \(4/3\) is the optimal threshold for each. The supplied excerpt states that replacing the bounds by \(\alpha>\alpha_0\) with \(\alpha_0<4/3\) does not give valid conclusions in general. It does not specify the counterexamples or the behavior at \(\alpha=4/3\). The admissible graph and rate quantifiers, the probability qualifier for the conclusions, and the precise meaning of optimality require recovery from the complete original problem statements. Consequently, this is a conditional reconstruction of the target, not a complete specification of the stochastic model.

The supplied original-statement field contains the labels ``Open Problem 1'' and ``Open Problem 2'' without their statements, followed by the lower-bound observation and a separate discussion of the critical case \(\alpha=1\). That discussion says that a function \(L\), not defined in the record, is non-strictly concave and that equilibrium need not be unique. For an even cycle of length \(2k\) with \(p(v)\equiv1\), it gives equilibria
\[
\{(a,b,a,b,\ldots,a,b):a+b=2\},
\]
without specifying the admissible domain of \(a,b\). It states that odd cycles have a unique equilibrium and that, as the even-cycle length increases, the law of \((a,b)\) converges to the Dirac measure at \((1,1)\); the equilibrium selection and probability law are not defined. Motivated by viewing \(\Z\) as an infinite cycle, the excerpt also records a separate Conjecture 2: for \(G=\Z\) and \(p(v)\equiv1\), almost surely \(x_v(t)\to1\) as \(t\to\infty\) for every \(v\). The quantities \(x_v(t)\) and their normalization are not defined here. These critical-case statements are supplied context and are not substituted for the supercritical threshold problem.
\subsection{Short English statement}
Determine whether the surviving components are trees for every \(\alpha>4/3\), and whiskers for every \(\alpha>4/3\) when all firing rates equal one, and whether these thresholds are optimal. The supplied record omits the stochastic model, graph hypotheses, probability qualifier, and definition of a whisker, so a fully quantified statement requires source review.
\subsection{Sources}
\begin{itemize}
\item Local UnsolvedMath registry, supplied record with permanent ID \(30005455\), source code \texttt{catalogue entry}, title ``Sharp Reinforcement Thresholds for Tree and Whisker Components.'' The record belongs to the Oberwolfach Reports Open Problems collection and gives the year 2023. No individual proposer is supplied.
\item \emph{MATRIX-MFO Tandem Workshop: Stochastic Reinforcement Processes and Graphs} (2023), \emph{Oberwolfach Reports}, pp. 639--679, DOI 10.4171/owr/2023/12: \url{https://doi.org/10.4171/owr/2023/12}. This is the supplied source report and is also listed as the ``original Oberwolfach report'' among the literature sources.
\item Hirsch--Holmes--Kleptsyn, \emph{Infinite WARM graphs III}: \url{https://doi.org/10.1088/1361-6544/acc9a0}. This reference is supplied in the literature-source list; its contents have not been independently checked here.
\item Literature assessment dated exactly \texttt{2026-08-22}: the supplied curation labels the problem ``partially solved'' and reports finite surviving components, tree or whisker structure in sufficiently large-\(\alpha\) regimes, and phase transitions, while saying that the sharp lower bounds requested by the two problems remain unresolved. This is supplied dated curation, not an independently verified literature assessment; its status and cited claims require independent verification.
\item The registry labels the statement \texttt{corrected\_verified} and says it was checked against the official report and preserves its hypotheses, notation, and target. This is a supplied verification claim, not verification performed for this draft. The original excerpt's reference label [8] has no matching bibliographic entry in the supplied record.
\end{itemize}
\end{document}
